% !TeX root = ../papers/nc-p3.tex
\section{小さいPlücker射影中心と共通三次因子曲線の次数}
\label{sec:deficient-plucker-centers}
\begin{theorem}\label{thm:deficient-plucker-centers}
特異核の八余因子有限鉛筆元で、定理\ref{thm:singular-quintic-pfaffians}の十五個の四次式が張る次元は11以下である。
その階数3以下軌跡の一次元部分の次数は、重複度を含めて少なくとも4である。
その各既約曲線の台は入力直線または滑らかな平面二次曲線である。
\end{theorem}
\begin{proof}[過剰曲線の次数下界]
四次式の包の次元を $r$、最大小行列式イデアルを $J$、その零スキームを $Y$ とする。
導手直線以外の $Y$ の各曲線は階数1軌跡に入り、一般の非比例点では固有直線入射に一意に戻る。
この入射はスカラー無限遠を避けるため射影曲線上でスカラーは一定であり、
既出のスカラー平面式から入力曲線は平面に入る。比例点は導手直線上にしかない。
次数3以上の整平面曲線が階数1軌跡に入ることはない。その二次Veroneseの線型包はベクトル次元6であり、
出力の曲線が非退化な四切断を持てばその平方積は少なくとも七次元である（一般滑らかな点で異なる先頭位数の基底を取り、
最小・最大位数との積を数える）。従って像も平面に入り、全二次式の座標変換が入力平面の平方六空間を
出力平面の平方六空間に送る。入力平面全体で $N$ が階数3以下となり、整な $D$ に反する。
これで台の主張と $Y$ の純次元が高々1であることを得る（孤立点は許す）。

四変数環上の六行四列Eagon--Northcott複体を取る。$J$ の高さは2で、唯一の正ホモロジーの可能性は $H_1$ である。
実際最終二微分は期待階数を $V(J)$ の外で持つから、この二写像からなる長さ2の自由複体は
Buchsbaum--Eisenbud判定で正次数に完全である。その終端余核の深さは2以上であり、そこから
無ねじれの $J$ への射の核 $H_1$ も深さ2以上である。$H_1$ は $Y$ の錐に支えられる。
Euler分子は $1-15t^4+24t^5-10t^6$ で対応するHilbert多項式は定数20なので、
$H_1$ は二次元Cohen--Macaulayで多重度は $Y$ の曲線次数に等しい。
次数4がその最初の可能な次数で、そこで $15-r$ 個の独立な定数関係がある。
一般の二一次式のArtin還元でこの次数は減らず、その全長は多重度であるから
\[
                            \deg Y\ \ge\ 15-r .
\]
\end{proof}

\begin{proof}[中心直線・中心平面の排除]
$B=\operatorname{Bl}_X\mathbb P^3$ は指定辺の四次式系によって $\operatorname{Gr}(2,R)\subset\mathbb P^{14}$ に埋め込まれる。
定理\ref{thm:octic-kernel-canonical-model}の直線切断による埋込み証明は任意のこの有限辺に通用する。
$J$ の生成空間の中心を $\Lambda\simeq\mathbb P^{14-r}$ とすると、
$X$ の外で $Y$ は $B\cap\Lambda$ の入力像にスキームとして一致する。
$B$ の非例外直線からの入力像は直線である。Grassmann直線の全二空間は固定関係 $u\in R$ を含み、
入力は $u$ の左核に入る。純関係がないのでその射影核は高々直線である。
$B$ 内の平面は例外平面に限られる。入力束が平面上 $\mathcal O(d)$、$d>0$ なら射は有限であり、
例外因子とこの平面の交叉は有限点へ縮約されるから因子としては空である。
するとPlücker束 $4H-E$ の次数1が $4d$ に等しくなってしまう。
よって $r=14$ の点中心では曲線がなく、$r=13$ の直線中心では入力次数は高々1で下界に反する。

$r=12$ の平面中心がGrassmann多様体に含まれなければ、その交叉は一つの非零二次式で切られる
平面二次曲線の部分スキームである。直線成分だけなら重複度込みで入力次数は高々2である。
滑らかな二次曲線ならその上の普遍二空間束は $\mathcal O(-2)\oplus\mathcal O$ または $\mathcal O(-1)^2$ となる。
前者は固定関係を含んで入力像は直線である。後者は二本の一次列による計八本の一次式係数の条件で
入力四座標の共通核を与える。核は一点であり、一般に独立な三行の三次小行列式から入力次数は3以下である。
次数3なら四つの三次小行列式には共通零点がないので、この三行は
$\mathcal O^4\twoheadrightarrow\mathcal O(1)^3$ を与え、双対完全列で全 $H^0(\mathcal O(3))$ が入力座標に現れる。
その像はねじれ三次曲線となるが、$Y$ にこの曲線は存在しない。従って次数3の下界に達しない。

中心平面全体がGrassmann多様体に入るなら、全二空間が固定関係を含む型か、固定三空間 $T\subset R$ に含まれる型である。
前者はスキームとしても入力が固定関係の左核に入り、次数は高々1である。後者を扱う。
入力 $x$ による $T$ の評価を $A_T(x):T\to k^4$ とする。これが零の入力点があれば、その点では
六関係評価の核は正確に三空間 $T$ となり、中心平面はそこの例外平面そのものである。
それ以外では $x\mapsto A_T(x)$ は線型埋込みで、中心切断上で $A_T$ は随所で階数1である。
入力曲線の次数は3以下である。実際一般入力平面でその次数はSegre多様体の有限な平面切断長以下である。
入力の $X$ を避ければ $A_T$ の階数1条件はまさに $A(x)$ の二次元核が $T$ に入る条件で、
スキームとしても一致する（単元小行列式で核束を復元する）。

ここでSegre多様体の任意の有限な平面切断は長さ3以下である。二次式で切られるから最大4であり、
4なら $(2,2)$ 完全交叉で全平面を張る。両因子の一次評価の次元を $u,v$ とすると、積は平面の一次式であり、
$\omega=\mathcal O(1)$ の完全跡に消されるので $u+v\le4$ である。一方が1なら切断の線型包が
一方を固定するSegre線型空間に入って全平面が残る。従って $u=v=2$ であるが、今度は全平面が
$\mathbb P^1\times\mathbb P^1$ の三次元線型包に入り、二次曲線が残ってしまう。

中心平面内の切断の一次元部分を平面因子 $C_0$ とする（各成分の一般重複度で取る）。
これは切断の閉部分であり、他の例外平面とは高々一点で交わる。実際二つの異なる三空間に含まれる二空間は高々一つである。
$A_T$ のSegre二射影の第一射影は中心座標そのもの、第二射影の線束を $J_0$ とすると
入力の超平面束は $\mathcal O_{C_0}(1)\otimes J_0$ となる。
従って $\deg C_0+\deg J_0=\deg Y\le3$ である。
$\deg C_0=3$ なら $J_0$ の射は被約台の各成分で一定、平面因子の連結性で同一点に行く。
$H^0(\mathcal O_{C_0})=k$ なのでアフィンチャートへのこの射は厚みを含めて一定である。
すると $A_T$ の線型像が固定第二因子との積の平面を含む（$C_0$ の線型包は全中心平面）。
階数1入力に平面が生じて矛盾する。従って $\deg C_0\le2$ であり、
直線または二次曲線についての先の次数評価が再び使える。これで $r=12$ も排除された。
\end{proof}

\path{computations/verify_deficient_plucker_centers.py} は完全交叉外でのEuler分子と、
次数3に達する二次曲線の核が入力の四つの独立な三次式を与える計算を確認する。
